% s02-3-reduction.tex —— 2.3 节：D-mod(Bun^I_G) 的紧生成（一）：归约与全局生成（印刷页 36–44）

\section{\texorpdfstring{$D\text{-}\mathrm{mod}(\Bun^I_G)$ 的紧生成}{D-mod(IG) 的紧生成}}
\label{sec:bunI-cg}

本节证明本文的主结果，即 $D\text{-}\mathrm{mod}(\Bun^I_G)$ 紧生成。

\subsection{归约步骤}
\label{sec:reduction}

与非分歧情形一样，策略是把 $\Bun^I_G$ 写成拟紧余截断开子叠之并，从而说明
$\Bun^I_G$ 可截断。由于 $\Bun^I_G$ 局部 QCA，随后可引用命题~\ref{prop:truncatable}
断定 $D\text{-}\mathrm{mod}(\Bun^I_G)$ 紧生成。

自始至终固定一点 $x\in X$（Iwahori 层级结构所在处），它诱导映射
\[
\Bun_G=\Maps(X,\pt/G)\to \Maps(\pt,\pt/G)=\pt/G.
\]
利用此映射，定义

\begin{definition}
\label{def:bunI}
$X$ 上带 $x\in X$ 处 Iwahori 层级结构的主 $G$-丛模叠，由纤维积给出：
\[
\Bun^I_G:=\Bun_G\times_{\pt/G}\pt/B.
\]
\end{definition}

注意有忘掉 Iwahori 层级结构的典范映射
\[
\Bun^I_G\to \Bun_G.
\]
对 $\lambda\in \Lambda^{+,Q}_G$，记
\[
\Bun^{I,(\lambda)}_G:=\Bun^{(\lambda)}_G\times_{\pt/G}\pt/B
=\Bun^{(\lambda)}_G\times_{\Bun_G}\Bun^I_G,
\]
\srcpage{36}
对 $S\subset \Lambda^{+,Q}_G$ 记
\[
\Bun^{I,(S)}_G:=\Bun^{(S)}_G\times_{\pt/G}\pt/B
=\Bun^{(S)}_G\times_{\Bun_G}\Bun^I_G .
\]
由于映射
\[
\Bun^I_G\to \Bun_G
\]
有限型，开子叠
\[
\Bun^{I,(\le\theta)}_G:=\Bun^{(\le\theta)}_G\times_{\Bun_G}\Bun^I_G
\]
拟紧，并且
\[
\Bun^I_G=\bigcup_{\theta\in \Lambda^{+,Q}_G}\Bun^{I,(\le\theta)}_G .
\]
因此，与非 Iwahori 层级结构的情形一样，只需证

\begin{proposition}
\label{prop:bunI-cotrunc}
只要对一切单根 $\check\alpha_i\in \Gamma_G$ 都有
\[
\langle\theta,\check\alpha_i\rangle\ge 2g-1,
\]
$\Bun^{I,(\le\theta)}_G$ 就是余截断的。
\end{proposition}

注意界已从 $2g-2$ 变为 $2g-1$。原因后文自会清楚。先作定义：

\begin{definition}
\label{def:iwahori-padmissible}
称 $P$-可容子集 $S\subset \Lambda^{+,Q}_G$ 是\emph{Iwahori-收缩 $P$-可容的}
（Iwahori-contractively $P$-admissible），若 $S$ 是 $P$-可容的，且对一切
$\lambda\in S$，只要 $i\notin \Gamma_M$，就有
\[
\langle\lambda,\check\alpha_i\rangle>2g-1.
\]
\end{definition}
\srcpage{37}
为说明 $\Bun^{I,(\le\theta)}_G$ 余截断，注意对 $\theta'\ge \theta$，如 2.1.6 小节那样，
可把
\[
\Bun^{I,(\le\theta')}_G\setminus \Bun^{I,(\le\theta)}_G
\]
写成有限个 $\Bun^{I,(S)}_G$ 之并，其中 $S\subset\Lambda^{+,Q}_G$ 是 Iwahori-收缩
$P$-可容的（$P$ 可变动）。

现在到了与非分歧情形证明的第一处分歧：我们将按\emph{相对位置}把上述子叠
$\Bun^{I,(S)}_G$ 进一步分层。

设 $S$ 如上为 Iwahori-收缩 $P$-可容。如同定义 $\Bun^I_G$，定义
\[
\Bun^I_P:=\Bun_P\times_{\pt/G}\pt/B,
\]
并类似地记
\[
\Bun^{I,(S)}_P:=\Bun^{(S)}_P\times_{\pt/G}\pt/B.
\]
注意有由 $\Bun^{(S)}_P\to \Bun_G$ 基变换得到的映射
\[
\Bun^{I,(S)}_P\to \Bun^{I}_G,
\]
它诱导同构
\[
\Bun^{I,(S)}_P\xrightarrow{\ \sim\ }\Bun^{I,(S)}_G .
\]
叠
\[
\pt/P\times_{\pt/G}\pt/B=P\backslash G/B
\]
被局部闭子叠
\[
P\backslash PwB/B\subset P\backslash G/B
\]
\srcpage{38}
分层，其中 $w$ 遍历 $W_M\backslash W$。映射
\[
\Bun^{I,(S)}_P\to \Bun^I_P\to \pt/P\times_{\pt/G}\pt/B
\]
使我们得以定义

\begin{definition}
\label{def:w-piece}
$\Bun^{I,(S)}_P$ 中相对位置为 $w$ 的片（piece），由纤维积给出：
\[
{}^{w}\Bun^{I,(S)}_P
=\Bun^{I,(S)}_P\times_{P\backslash G/B}P\backslash PwB/B .
\]
利用同构 $\Bun^{I,(S)}_P\xrightarrow{\sim}\Bun^{I,(S)}_G$，同时把
${}^{w}\Bun^{I,(S)}_G$ 定义为 ${}^{w}\Bun^{I,(S)}_P$ 的像。
\end{definition}

由于有限个截断局部闭子叠之并截断，命题~\ref{prop:bunI-cotrunc} 归约为

\begin{proposition}
\label{prop:contraction-target}
局部闭子叠
\[
{}^{w}\Bun^{I,(S)}_G\to \Bun^I_G
\]
收缩，因而截断。
\end{proposition}

\subsection{概览}
\label{sec:outline}

本节余下部分致力于证明命题~\ref{prop:contraction-target}，即：对 Iwahori-收缩
$P$-可容子集 $S$，${}^{w}\Bun^{I,(S)}_G\to \Bun^I_G$ 收缩。

记 $\uB=B\cap M$，并定义
\[
\Bun^{I_M}_M:=\Bun_M\times_{\pt/M}\pt/\uB,
\]
并类似地
\[
\Bun^{I_M,(S)}_M:=\Bun^{(S)}_M\times_{\pt/M}\pt/\uB .
\]
\srcpage{39}
对 $w\in W_M\backslash W$，在 $W$ 中有唯一的最短长度提升，今后也（滥用地）记作 $w$。
对这种 $w$，群同态 $\uB\to G$（$b\mapsto wbw^{-1}$）的像含于 $B$，故得到叠的映射
\[
\pt/\uB\xrightarrow{\ w\ }\pt/B.
\]
图
\begin{center}
\begin{tikzpicture}[>=stealth,scale=1.0]
\node (a) at (0,0) {$\pt/\uB$};
\node (b) at (-2.2,-1.1) {$\pt/M$};
\node (c) at (2.2,-1.1) {$\pt/B$};
\node (d) at (-4.0,-2.2) {$\pt/P$};
\node (e) at (0.9,-2.9) {$\pt/G$};
\draw[->] (a)--(b);
\draw[->] (a)--node[pos=0.75,above right]{$w$}(c);
\draw[->] (b)--(d);
\draw[->] (c)--(e);
\draw[->] (d)--(e);
\end{tikzpicture}
\end{center}
交换，并且
\[
\pt/\uB\to \pt/P\times_{\pt/G}\pt/B
\]
经子叠 $P\backslash PwB/B\subset P\backslash G/B$ 分解。于是图
\srcpage{40}
\begin{center}
\begin{tikzpicture}[>=stealth,scale=0.95]
\node (t) at (0,0) {$\Bun^{I_M,(S)}_M$};
\node (m) at (-3.4,-1.2) {$\Bun^{(S)}_M$};
\node (b) at (3.4,-1.2) {$\pt/\uB$};
\node (p) at (-6.6,-2.4) {$\Bun^{(S)}_P$};
\node (mm) at (0,-2.4) {$\pt/M$};
\node (bb) at (6.6,-2.4) {$\pt/B$};
\node (pp) at (-3.4,-3.6) {$\pt/P$};
\node (g) at (0,-4.8) {$\pt/G$};
\draw[->] (t)--(m);
\draw[->] (t)--(b);
\draw[->] (m)--(p);
\draw[->] (m)--(mm);
\draw[->] (b)--(mm);
\draw[->] (b)--node[pos=0.8,above right]{$w$}(bb);
\draw[->] (p)--(pp);
\draw[->] (mm)--(pp);
\draw[->] (bb)--(g);
\draw[->] (pp)--(g);
\end{tikzpicture}
\end{center}
诱导出映射
\[
\Bun^{I_M,(S)}_M\to \Bun^{I,(S)}_P,
\]
它经 ${}^{w}\Bun^{I,(S)}_P\subset \Bun^{I,(S)}_P$ 分解。

完全类似地，图
\[
\begin{tikzpicture}[>=stealth,scale=1.0]
\node (a) at (0,0) {$\pt/\uB$};
\node (b) at (-2.2,-1.1) {$\pt/M$};
\node (c) at (2.2,-1.1) {$\pt/B$};
\node (d) at (-4.0,-2.2) {$\pt/P^-$};
\node (e) at (0.9,-2.9) {$\pt/G$};
\draw[->] (a)--(b);
\draw[->] (a)--node[pos=0.75,above right]{$w$}(c);
\draw[->] (b)--(d);
\draw[->] (c)--(e);
\draw[->] (d)--(e);
\end{tikzpicture}
\]
\srcpage{41}
交换，并且 $\pt/\uB\to \pt/P^-\times_{\pt/G}\pt/B$ 经子叠
$P^-\backslash P^-wB/B\subset P^-\backslash G/B$ 分解。于是相应的图诱导出映射
\[
\Bun^{I_M,(S)}_M\to \Bun^{I,(S)}_{P^-},
\]
它经 ${}^{w}\Bun^{I,(S)}_{P^-}\subset \Bun^{I,(S)}_{P^-}$ 分解。

于是有交换图
\[
\begin{tikzcd}[column sep=huge,row sep=large]
\Bun^{I_M,(S)}_M \arrow[r] \arrow[d] & {}^{w}\Bun^{I,(S)}_{P^-} \arrow[d]\\
{}^{w}\Bun^{I,(S)}_P \arrow[r] & \Bun^I_G .
\end{tikzcd}
\]
结合 2.1.4 小节，命题~\ref{prop:contraction-target} 归约为证明

\begin{proposition}
\label{prop:four-conditions}
在上述图中，
\begin{enumerate}
\item[(a)] $\Bun^{I_M,(S)}_M\to {}^{w}\Bun^{I,(S)}_P\times_{\Bun^I_G}{}^{w}\Bun^{I,(S)}_{P^-}$ 是开嵌入；
\item[(b)] $\Bun^{I_M,(S)}_M\to {}^{w}\Bun^{I,(S)}_P$ 满；
\item[(c)] ${}^{w}\Bun^{I,(S)}_{P^-}\to \Bun^I_G$ 光滑；
\item[(d)] $\Bun^{I_M,(S)}_M\to {}^{w}\Bun^{I,(S)}_{P^-}$ 收缩。
\end{enumerate}
\end{proposition}

命题~\ref{prop:four-conditions} 的 (a) 将在 2.3.4 小节证明，(b) 在 2.3.5 小节，
(c) 在 2.3.6 小节，(d) 在 2.3.7 小节。

\subsection{\texorpdfstring{$n(P)_{F_M}$ 的全局生成}{n(P)FM 的全局生成}}
\label{sec:nP-global}

本小节说明：在 $F_M\in \Bun_M$ 上附加一定条件时，$n(P)_{F_M}$ 全局生成。先回顾全局生成的概念。

\begin{definition}
\label{def:globally-generated}
称 $X$ 上的向量丛 $E$ 是\emph{全局生成的}，若对一切 $x\in X$，映射
\[
H^0(X,E)\to E_x
\]
满。
\end{definition}
\srcpage{42}
以下是检验曲线上向量丛全局生成的一些有用事实。

\begin{lemma}
\label{lem:gg-h1}
向量丛 $E$ 全局生成，当且仅当对一切 $x\in X$ 有 $H^1(X,E(-x))=0$。
\end{lemma}

\begin{lemma}
\label{lem:gg-filtration}
设 $E$ 带滤过
\[
0=E_0\subset E_1\subset \cdots \subset E_m=E,
\]
其子商 $E_i/E_{i-1}$ 都是向量丛，并设
\[
\gr(E)=\bigoplus_i E_i/E_{i-1}
\]
为该滤过下 $E$ 的伴随分次。若 $H^1(X,\gr(E))=0$，则 $H^1(X,E)=0$。
\end{lemma}

由前两个引理立即得到

\begin{lemma}
\label{lem:gg-gr}
设 $E$ 如上。若 $H^1(X,\gr(E))=0$ 且 $\gr(E)$ 全局生成，则 $E$ 全局生成。
\end{lemma}

本小节余下部分用上述关于全局生成的基本事实说明 $n(P)_{F_M}$ 全局生成。

\begin{lemma}
\label{lem:V-global}
设 $\tilde G$ 为约化群，$V$ 为有限维 $\tilde G$-表示，$Z_0(\tilde G)$ 经 $\check\mu$
作用其上，$F_{\tilde G}\in \Bun^{\lambda,ss}_{\tilde G}$。若
\[
\langle\lambda,\check\mu\rangle>2g-1,
\]
则 $V_{F_{\tilde G}}$ 全局生成。
\end{lemma}

\begin{proof}
由 Serre 对偶，
\[
H^1(X,V_{F_{\tilde G}}(-x))
=H^0\bigl(X,(V^*)^{F_{\tilde G}}(x)\otimes \omega_X\bigr)^*.
\]
由 \cite[引理 10.3.2]{DG15}，$(V^*)^{F_{\tilde G}}$ 的每个线子丛次数
$\le -\langle\lambda,\check\mu\rangle<-(2g-1)$，故
$(V^*)^{F_{\tilde G}}(x)\otimes\omega_X$ 的每个线子丛次数 $<0$。因此
\[
H^0\bigl(X,(V^*)^{F_{\tilde G}}(x)\otimes\omega_X\bigr)=0,
\]
再由引理~\ref{lem:gg-h1} 即得结论。
\end{proof}
\srcpage{43}

\begin{proposition}
\label{prop:nP-global}
设 $F_M\in \Bun^{(\lambda)}_M$ 满足：只要 $i\notin \Gamma_M$ 就有
$\langle\lambda,\check\alpha_i\rangle>2g-1$。则相伴向量丛 $n(P)_{F_M}$ 的 $H^1$ 为 $0$，
且全局生成。
\end{proposition}

\begin{proof}
仿 \cite[命题 10.4.5]{DG15} 的证明：令 $P_\lambda$ 为 $M$ 中对应子集
\[
\{i\in \Gamma_M\mid \langle\lambda,\check\alpha_i\rangle=0\}\subset \Gamma_M
\]
的抛物子群，$M_\lambda$ 为其 Levi。由于 $F_M\in \Bun^{(\lambda)}_M$，由
\cite[定理 7.4.3]{DG15}，$F_M$ 有次数为 $\lambda$ 的、到 $P_\lambda$ 的半稳定归约，
记作 $F_{P_\lambda}$。记 $F_{M_\lambda}$ 为沿 $P_\lambda\to M_\lambda$ 诱导的主
$M_\lambda$-丛。则 $n(P)_{F_{M_\lambda}}$ 是 $n(P)_{F_M}=n(P)_{F_{P_\lambda}}$ 的
一个伴随分次。该处已证 $H^1(X,n(P)_{F_{M_\lambda}})=0$，故由引理~\ref{lem:gg-filtration}
得 $H^1(X,n(P)_{F_M})=0$。

至于全局生成，由引理~\ref{lem:gg-gr}，只需说明 $n(P)_{F_{M_\lambda}}$ 全局生成。
现注意（如 \cite[10.4.1 小节]{DG15}），$n(P)$ 可写成 $M_\lambda$-表示的直和
\[
\bigoplus V_{M_\lambda,\check\alpha},
\]
其中直和遍历 $G$ 的正根 $\check\alpha$，使得存在某个 $i\notin \Gamma_M$ 满足
\srcpage{44}
$\mathrm{coeff}_i(\check\alpha)>0$。$Z_0(M_\lambda)$ 经 $\check\alpha$ 作用在
$V_{M_\lambda,\check\alpha}$ 上，且由
\[
\langle\lambda,\check\alpha\rangle
\ge \langle\lambda,\check\alpha_i\rangle\cdot \mathrm{coeff}_i(\check\alpha)
\ge \langle\lambda,\check\alpha_i\rangle\cdot 1>2g-1,
\]
引理~\ref{lem:V-global} 给出 $V_{M_\lambda,\check\alpha,F_{M_\lambda}}$ 全局生成。
\end{proof}

\begin{corollary}
\label{cor:UP-surj}
设 $F_M$ 如命题~\ref{prop:nP-global}。则映射
\[
H^0(X,U(P)_{F_M})\to U(P)_{x^*F_M}\xrightarrow{\ \sim\ }U(P)
\]
满。
\end{corollary}

\begin{proof}
特征 $0$ 下存在双射的 $M$-等变映射 $n(P)\to U(P)$。
\end{proof}

\begin{corollary}
\label{cor:gp-global}
设 $F_{P^-}\in \Bun^{(\lambda)}_{P^-}$ 满足：只要 $i\notin \Gamma_M$ 就有
$\langle\lambda,\check\alpha_i\rangle>2g-1$。则 $(\mathfrak{g}/\mathfrak{p}^-)_{F_{P^-}}$
全局生成。
\end{corollary}

\begin{proof}
令 $F_M$ 为沿映射 $P^-\to M$ 诱导的主 $M$-丛。则
$(\mathfrak{g}/\mathfrak{p}^-)_{F_M}=n(P)_{F_M}$ 是
$(\mathfrak{g}/\mathfrak{p}^-)_{F_{P^-}}$ 的伴随分次。又 $H^1(X,n(P)_{F_M})=0$，
故结论由引理~\ref{lem:gg-gr} 得出。
\end{proof}
